\lecture{1.4}{Independence and Conditional Independence}
\label{lec:1-4}

\subsection{Definition of Independence}

Independence formalizes the idea that learning one piece of information does not change what we know about another.

\begin{definition}[Independence]
Events $B$ and $C$ are independent if
\[
\Prob(B,C)=\Prob(B)\Prob(C).
\]
Equivalently, when $\Prob(C)>0$ and $\Prob(B)>0$,
\[
\Prob(B\given C)=\Prob(B)
\qquad\text{and}\qquad
\Prob(C\given B)=\Prob(C).
\]
For random variables $X$ and $Y$, we write $X\ci Y$ if all events determined by $X$ are independent of all events determined by $Y$.
\end{definition}

The equivalence follows directly from conditional probability. If $\Prob(B,C)=\Prob(B)\Prob(C)$, then
\[
\Prob(B\given C)=\frac{\Prob(B,C)}{\Prob(C)}
=\frac{\Prob(B)\Prob(C)}{\Prob(C)}
=\Prob(B).
\]
Thus conditioning on $C$ does not change the probability of $B$.

\begin{keyidea}
Independence means that one variable carries no information about the other in the population.
\end{keyidea}

\subsection{Definition of Conditional Independence}

Conditional independence is independence within strata of another variable.

\begin{definition}[Conditional independence]
Events $B$ and $C$ are conditionally independent given $D$ if
\[
\Prob(B,C\given D)=\Prob(B\given D)\Prob(C\given D)
\]
whenever the conditioning probabilities are well-defined. Equivalently,
\[
\Prob(B\given C,D)=\Prob(B\given D).
\]
For random variables, we write $X\ci Y\given Z$.
\end{definition}

Plainly, $X\ci Y\given Z$ means that after restricting to people with the same value of $Z$, knowing $Y$ gives no additional information about $X$.

Conditional independence is not the same as marginal independence. It is possible for two variables to be associated in the full population but independent within strata of a third variable. It is also possible for two variables to be independent in the full population but associated after conditioning.

\subsection{Independence in Randomized Trials}

Consider a randomized trial in which $A$ indicates assignment to a vaccine arm, with
\[
\Prob(A=1)=\frac{1}{2}
\]
for every participant. Let $W$ denote age. Because the treatment assignment is generated by a random device that ignores age,
\[
\Prob(A=1\given W=w)=\Prob(A=1)=\frac{1}{2}
\]
for every age $w$ in the trial population. Hence $A\ci W$.

\begin{assumptionbox}
In an individually randomized trial with equal allocation, treatment assignment is independent of baseline covariates:
\[
A\ci W.
\]
In plain language, baseline characteristics do not help predict treatment assignment because the randomization mechanism did not use them.
\end{assumptionbox}

This independence is a design property. It is not discovered by observing that the treatment groups look similar in one finite sample. In a small randomized trial the realized treatment group may contain older participants by chance, but the assignment mechanism is still independent of age.

\subsection{Stratified Randomization}

Now suppose the trial uses stratified randomization by age. Participants younger than 65 receive vaccine with probability $1/2$, while participants aged at least 65 receive vaccine with probability $2/3$:
\[
\Prob(A=1\given W<65)=\frac{1}{2},
\qquad
\Prob(A=1\given W\ge 65)=\frac{2}{3}.
\]
In this design, $A$ and $W$ are not marginally independent. Knowing that someone is at least 65 changes the probability of treatment assignment.

However, within age strata, the assignment may still be randomized with respect to other baseline covariates. Let $Z$ denote congestive heart failure status. If the randomization rule uses age but not $Z$, then
\[
A\ci Z\given W.
\]
Among people of the same age stratum, knowing congestive heart failure status gives no additional information about treatment assignment.

\begin{figure}[H]
\centering
\begin{tikzpicture}[node distance=2.2cm,>=Latex]
  \node[draw, circle, fill=blue!5] (W) {$W$\\age};
  \node[draw, circle, fill=green!5, right=of W] (A) {$A$\\treatment};
  \node[draw, circle, fill=orange!6, below=of W] (Z) {$Z$\\CHF};
  \draw[->, thick] (W) -- (A);
  \draw[<->, thick, dashed] (W) -- (Z);
  \node[align=center, right=1.1cm of A] {Stratified randomization:\\age affects assignment probability.};
\end{tikzpicture}
\caption{In a stratified trial, treatment assignment can depend on the stratifying variable while remaining independent of other covariates within strata.}
\end{figure}

\subsection{Why Conditioning Can Remove Information}

Conditioning can remove information because it compares units only within a fixed subgroup. In the stratified randomization example, congestive heart failure status may be associated with age. In the full population, learning $Z$ tells us something about likely age, and age tells us something about treatment probability. Therefore $Z$ can be informative about $A$ marginally.

Once we condition on age, that indirect information path is removed. Among participants in the same age stratum, $Z$ no longer conveys information about the randomization probability. This is why
\[
A\nci Z
\qquad\text{but}\qquad
A\ci Z\given W
\]
can both be true.

The following table gives a concrete illustration:
\[
\begin{array}{c|cc}
& \Prob(A=1\given W,Z=0) & \Prob(A=1\given W,Z=1)\\
\hline
W<65 & 1/2 & 1/2\\
W\ge 65 & 2/3 & 2/3
\end{array}
\]
Within each age stratum, the treatment probability does not depend on $Z$. If people with $Z=1$ are more often in the older stratum, then $\Prob(A=1\given Z=1)$ can still exceed $\Prob(A=1\given Z=0)$ marginally.

\subsection{Why Conditional Independence Matters for Causal Inference}

Causal inference often requires comparing treated and untreated people who are similar with respect to variables that affect treatment and outcome. The mathematical language for this idea is conditional independence.

In later chapters, a central causal assumption for a binary treatment will take the form
\[
\Ya{a}\ci A\given W,
\qquad a\in\{0,1\}.
\]
This says that, among people with the same covariates $W$, treatment assignment $A$ carries no additional information about the counterfactual outcome under treatment level $a$. This is a causal version of "as-if randomized within strata of $W$."

The probability concepts in this lecture are therefore not merely technical preliminaries. They are the language used to state whether observed data can support causal comparisons.

\begin{keyidea}
Marginal independence asks whether two variables are unrelated in the whole population. Conditional independence asks whether they are unrelated after comparing only units in the same stratum of the conditioning variables.
\end{keyidea}
